% Abstract, bound from editorial/upstream/front_matter_rewrite.tex at the
% Session 37 Stage 4 front-matter bind (author-approved, D303/D306) and
% trimmed to the 250-word gate at Session 38 per Carlos's instruction
% (2026-07-11): micro-trims only, no content removed.
% REVIEW-CLAIM: (Session 38) "the predictive coordinates are the most
% forecastable" corrected to "the wake-supervised states are the most
% forecastable": the D310 shared-operator recomputation makes the top block
% a statistical tie among the wake-headed states, so the draft's sentence
% (written pre-M1) is no longer supportable as stated.
Extreme vortex gusts drive a post-stall airfoil through a transient
reorganisation of its wake, and the release parameters alone do not
determine the flow after the vortex passes. We ask whether a reduced
state of the encounter can retain the wake dynamics, remain usable
under a low-order forecast model, and be estimated from the sparse wall
pressure a wing can carry. We compare predictive, reconstruction-trained
and proper-orthogonal-decomposition coefficient states at matched
dimension, on direct numerical simulations of a NACA 0012 airfoil at
$\alpha = 14^\circ$ and $\vRe = 5000$ perturbed by Taylor vortices
spanning gust ratio, core diameter and offset. A controlled supervision
study shows different ingredients supply the requirements.
Under the predictive objective considered, a scalar lift target is
necessary against latent collapse, a near-body force-element target
sharpens the peak loads, and a wake target is required for the wake
enstrophy and circulation; the multi-step predictive objective
contributes rollout stability rather than instantaneous readability.
Load information remains accessible at four coefficients, whereas
wake-field reconstruction requires at least sixteen. Under one
\DirectFC{} trained identically on every state, the wake-supervised
states are the most forecastable, autoregressive rollout degrades
through error accumulation, and supplying the true gust parameters
provides no benefit. A sequential ensemble filter sensing eight
wall-pressure taps then tracks the lift through the encounter, with
impact-phase $R^2 = \PoneTwoStageCovImpact$ (root-mean-square error
$\PoneTwoStageCovRmse$ in lift units) on the \TestSplit{} set and
$\PoneTwoStageCovImpactC$ at the held-out $|G| = 4$ boundary; the wake
enstrophy is not recovered by the wall-sensor configurations considered,
a limit of the sensing geometry rather than the state.
